% ============================================================
% HSBI Beamer — ISLP Lecture Series
% Mock Exam 2 — Solutions (Lectures 5-8, ISLP Chapters 4-6)
% Prof. Dr. Christoph Weisser — Summer Semester 2026
% Build: pdflatex -jobname=Mock_Exam_2_Solutions_Slides solutions_slides_2.tex  (twice)
% ============================================================
\documentclass[aspectratio=169,11pt]{beamer}
\usepackage[utf8]{inputenc}\usepackage[T1]{fontenc}\usepackage[english]{babel}
\usepackage{graphicx}\usepackage{booktabs}\usepackage{amsmath,amssymb,amsfonts}
\usepackage{mathtools}\usepackage{hyperref}\usepackage{xcolor}
\usepackage{verbatim}\usepackage{alltt}
\usepackage{tikz}\usetikzlibrary{calc,arrows.meta,positioning,shapes}
\usepackage{tcolorbox}\tcbuselibrary{skins}

\usetheme{Madrid}\usecolortheme{default}
\setbeamertemplate{navigation symbols}{}
\setbeamertemplate{footline}[frame number]
\setbeamertemplate{blocks}[rounded][shadow=false]
\setbeamertemplate{headline}{%
  \leavevmode\hbox{%
    \begin{beamercolorbox}[wd=\paperwidth,ht=2.2ex,dp=0.9ex]{section in head/foot}%
      {\fontsize{4.6}{5.2}\selectfont%
       \insertsectionnavigationhorizontal{\paperwidth}{}{\hskip0pt plus1filll}}%
    \end{beamercolorbox}}%
}
\setbeamerfont{section in head/foot}{size=\tiny}
\setbeamerfont{palette primary}{size=\tiny}
\setbeamerfont{palette tertiary}{size=\tiny}
\setbeamerfont{section in toc}{size=\footnotesize}

\usepackage{listings}
\definecolor{pyKw}{RGB}{0,0,180}\definecolor{pyStr}{RGB}{20,120,40}
\definecolor{pyCom}{RGB}{120,120,120}\definecolor{accent}{RGB}{38,70,140}
\lstdefinestyle{Pystyle}{language=Python,basicstyle=\ttfamily\footnotesize,
  keywordstyle=\color{pyKw}\bfseries,commentstyle=\color{pyCom}\itshape,
  stringstyle=\color{pyStr},showstringspaces=false,upquote=true,
  columns=fullflexible,breaklines=true,literate={~}{{$\sim$}}1}
\lstset{style=Pystyle}

\AtBeginSection[]{\begin{frame}{Outline}\footnotesize
  \tableofcontents[currentsection,sectionstyle=show/shaded,subsectionstyle=show/show/hide]
\end{frame}}

\newcommand{\islpsource}[1]{%
  \vfill\hfill{\tiny\itshape\color{gray}Source: #1 \textemdash{} ISLP, James et al.\ (2023).}\par
}

\newtcolorbox{takeaway}[1][Takeaway]{enhanced,colback=green!4,colframe=green!50!black,
  boxrule=0pt,leftrule=2.5pt,arc=2pt,fonttitle=\bfseries,title=#1,
  left=2mm,right=2mm,top=1mm,bottom=1mm}
\newtcolorbox{numexample}[1][Worked example]{enhanced,colback=orange!4,colframe=orange!70!black,
  boxrule=0pt,leftrule=2.5pt,arc=2pt,fonttitle=\bfseries,title=#1,
  left=2mm,right=2mm,top=1mm,bottom=1mm}
\newtcolorbox{readme}[1][How to read this]{enhanced,colback=blue!4,colframe=accent,
  boxrule=0pt,leftrule=2.5pt,arc=2pt,fonttitle=\bfseries,title=#1,
  left=2mm,right=2mm,top=1mm,bottom=1mm}

\title{Mock Exam 2 --- Solutions}
\subtitle{Lectures 5--8 (Chapters 4--6)}
\author{Prof.\ Dr.\ Christoph Weisser}\institute{Quantitative Research Methods \textperiodcentered\ HSBI}
\date{Summer Semester 2026}

\newtcolorbox{exercise}[1][Exercise]{enhanced, fontupper=\footnotesize, colback=purple!4, colframe=purple!55!black, boxrule=0pt, leftrule=2.5pt, arc=2pt, fonttitle=\bfseries, title=#1, left=2mm, right=2mm, top=1mm, bottom=1mm}
\newtcolorbox{solutionbox}[1][Solution]{enhanced, fontupper=\footnotesize, colback=teal!5, colframe=teal!55!black, boxrule=0pt, leftrule=2.5pt, arc=2pt, fonttitle=\bfseries, title=#1, left=2mm, right=2mm, top=1mm, bottom=1mm}
\newtcolorbox{longexercise}[1][Extended exercise (15 min)]{enhanced, fontupper=\footnotesize, colback=violet!6, colframe=violet!55!black, boxrule=0pt, leftrule=3pt, arc=2pt, fonttitle=\bfseries, title=#1, left=2mm, right=2mm, top=1mm, bottom=1mm}

% Reclaim ~2pt of body height (custom nav headline makes full slides tip over by ~1.83pt)
\addtobeamertemplate{frametitle}{}{\vspace{-2pt}}

\begin{document}
\begin{frame}[plain]\titlepage\end{frame}

% ------------------------------------------------------------
\begin{frame}{Overview --- problems and points}
  \footnotesize
  \begin{center}
  \begin{tabular}{@{}clc@{}}
    \toprule
    \textbf{Problem} & \textbf{Topic} & \textbf{Points}\\
    \midrule
    P1 & Logistic regression by hand & 18\\
    P2 & Discriminant analysis and naive Bayes --- concepts & 12\\
    P3 & Confusion matrix and ROC curve & 15\\
    P4 & Resampling methods (validation set $\cdot$ CV $\cdot$ bootstrap) & 18\\
    P5 & Model selection with $C_p$ and BIC & 12\\
    P6 & Ridge regression and the lasso & 15\\
    \midrule
     & \textbf{Total} \quad (90 minutes; points $\approx$ minutes) & \textbf{90}\\
    \bottomrule
  \end{tabular}
  \end{center}
  \vspace{1mm}
  \begin{takeaway}[How to use this deck]
    Attempt each problem under exam conditions before looking at the teal boxes.
    Allowed aids in the exam: a non-programmable calculator and one handwritten A4 sheet
    of notes (both sides). Round final results to 3 decimals.
  \end{takeaway}
\end{frame}

% ============================================================
% Problem 1
% ============================================================
\begin{frame}{Problem 1 --- restatement}
  \begin{exercise}[Problem 1 --- Logistic regression by hand (18 P)]
    A bank models the default probability of a credit-card customer as a function of the
    monthly balance $x$ (in EUR; observed range 0--400):
    \[
    \hat{p}(x)=\frac{e^{\hat{\beta}_0+\hat{\beta}_1x}}{1+e^{\hat{\beta}_0+\hat{\beta}_1x}},
    \qquad \hat{\beta}_0=-4,\quad \hat{\beta}_1=0.02 .
    \]
    \begin{itemize}
      \item[(a)] Predicted default probability at $x=100$ and $x=250$. \hfill (4 P)
      \item[(b)] Odds and log-odds at both balances; interpret the odds at $x=250$. \hfill (4 P)
      \item[(c)] Balance $x^{*}$ with $\hat{p}(x^{*})=0.5$. \hfill (3 P)
      \item[(d)] Factor change of the odds per EUR and per 100 EUR; why is
        ``probability $+0.02$ per EUR'' wrong? \hfill (4 P)
      \item[(e)] Two reasons why linear regression on the 0/1 response is inappropriate. \hfill (3 P)
    \end{itemize}
  \end{exercise}
\end{frame}

\begin{frame}{Problem 1 --- solution (a) and (b)}
  \begin{solutionbox}[Solution 1(a) --- predicted probabilities]
    Linear predictor (log-odds): $\hat{\eta}(x)=-4+0.02\,x$, so
    $\hat{\eta}(100)=-4+2=-2$ and $\hat{\eta}(250)=-4+5=1$. Hence
    \[
    \hat{p}(100)=\frac{e^{-2}}{1+e^{-2}}=\frac{0.1353}{1.1353}=\mathbf{0.119},
    \qquad
    \hat{p}(250)=\frac{e^{1}}{1+e^{1}}=\frac{2.7183}{3.7183}=\mathbf{0.731}.
    \]
  \end{solutionbox}
  \begin{solutionbox}[Solution 1(b) --- odds and log-odds]
    Odds $=\hat{p}/(1-\hat{p})=e^{\hat{\eta}(x)}$ and log-odds $=\hat{\eta}(x)$:
    \begin{center}
    \begin{tabular}{lcc}
      \toprule
       & $x=100$ & $x=250$\\
      \midrule
      log-odds & $-2.000$ & $1.000$\\
      odds & $e^{-2}=0.135$ & $e^{1}=2.718$\\
      \bottomrule
    \end{tabular}
    \end{center}
    Check with (a): $0.119/0.881=0.135$ and $0.731/0.269=2.718$. \checkmark\
    At $x=250$ the odds are about $2.72$ to 1: default is roughly $2.7$ times as probable
    as non-default.
  \end{solutionbox}
\end{frame}

\begin{frame}{Problem 1 --- solution (c) and (d)}
  \begin{solutionbox}[Solution 1(c) --- balance with a 50:50 prediction]
    $\hat{p}=0.5 \iff \text{odds}=1 \iff \text{log-odds}=0$:
    \[
    -4+0.02\,x^{*}=0 \quad\Longrightarrow\quad x^{*}=\frac{4}{0.02}=\mathbf{200\ \text{EUR}}.
    \]
    Below 200 EUR ``no default'' is predicted as more likely; above 200 EUR ``default''.
  \end{solutionbox}
  \begin{solutionbox}[Solution 1(d) --- interpretation on the odds scale]
    Per additional EUR the odds are \emph{multiplied} by $e^{\hat{\beta}_1}=e^{0.02}=1.020$
    ($+2.0\,\%$); per additional 100 EUR by $e^{100\cdot 0.02}=e^{2}=7.389$.
    The quoted statement is wrong because $\hat{\beta}_1$ acts on the \emph{log-odds}:
    the probability change per EUR is
    $\partial\hat{p}/\partial x=\hat{\beta}_1\,\hat{p}(x)\bigl(1-\hat{p}(x)\bigr)$,
    which depends on $x$ --- the S-curve is steepest at $\hat{p}=0.5$ and almost flat in
    the tails, so there is no constant probability effect.
  \end{solutionbox}
\end{frame}

\begin{frame}{Problem 1 --- solution (e)}
  \begin{solutionbox}[Solution 1(e) --- why not linear regression]
    \begin{itemize}
      \item[(1)] The straight line $\hat{\beta}_0+\hat{\beta}_1x$ is unbounded: for small or
        large balances it returns ``probabilities'' below 0 or above 1 --- not sensible
        estimates of $\Pr(Y=1\mid x)$. The logistic function always stays in $[0,1]$.
      \item[(2)] A 0/1 response violates the linear-model error assumptions: errors are
        non-normal and heteroscedastic --- $\operatorname{Var}(Y\mid x)=p(x)(1-p(x))$
        depends on $x$ --- so the usual inference is unreliable.
      \item[] (With $>2$ classes a numeric coding would additionally impose an artificial
        ordering and equal spacing between the classes.)
    \end{itemize}
  \end{solutionbox}
  \begin{takeaway}[Remember]
    Logistic regression is a \emph{linear model for the log-odds}:
    $\log\bigl(p/(1-p)\bigr)=\beta_0+\beta_1x$. Coefficients are additive on the log-odds
    scale and multiplicative on the odds scale --- never additive on the probability scale.
  \end{takeaway}
\end{frame}

% ============================================================
% Problem 2
% ============================================================
\begin{frame}{Problem 2 --- restatement}
  \begin{exercise}[Problem 2 --- Discriminant analysis and naive Bayes (12 P)]
    \begin{itemize}
      \item[(a)] LDA vs.\ QDA: the distinguishing covariance assumption; number of
        covariance matrices estimated with $K$ classes. \hfill (4 P)
      \item[(b)] Bias--variance: when does LDA beat QDA on test data and vice versa?
        Role of $n$ and of the true class covariances. \hfill (3 P)
      \item[(c)] Key assumption of naive Bayes; why can it help when $p$ is large relative
        to $n$? \hfill (3 P)
      \item[(d)] Which classifiers give linear and which give quadratic or more general
        boundaries? \hfill (2 P)
    \end{itemize}
  \end{exercise}
\end{frame}

\begin{frame}{Problem 2 --- solution (a) and (b)}
  \begin{solutionbox}[Solution 2(a) --- covariance assumptions]
    Both assume $X\mid Y=k\sim\mathcal{N}(\mu_k,\Sigma_k)$.
    \textbf{LDA:} one covariance matrix \emph{shared} by all classes
    ($\Sigma_1=\dots=\Sigma_K=\Sigma$) $\Rightarrow$ estimates \textbf{1} matrix
    ($p(p+1)/2$ parameters). \textbf{QDA:} a \emph{class-specific} $\Sigma_k$
    $\Rightarrow$ estimates \textbf{K} matrices ($K\,p(p+1)/2$ parameters).
  \end{solutionbox}
  \begin{solutionbox}[Solution 2(b) --- who wins when]
    QDA is far more flexible: \emph{low bias but high variance}; LDA is restrictive:
    \emph{low variance but biased if the $\Sigma_k$ truly differ}.
    \begin{itemize}
      \item \textbf{LDA wins} when $n$ is small (or $p$ large) --- variance reduction is
        crucial --- and when the class covariances are approximately equal (little bias
        incurred).
      \item \textbf{QDA wins} when $n$ is large --- its many parameters are estimated
        reliably --- and when the class covariances clearly differ, so the shared-$\Sigma$
        assumption of LDA would cause substantial bias.
    \end{itemize}
  \end{solutionbox}
\end{frame}

\begin{frame}{Problem 2 --- solution (c) and (d)}
  \begin{solutionbox}[Solution 2(c) --- naive Bayes]
    Assumption: \emph{within each class the $p$ predictors are independent}:
    $f_k(x)=f_{k1}(x_1)\cdots f_{kp}(x_p)$.
    Only $p$ one-dimensional marginals per class are estimated instead of a full
    $p$-dimensional joint density (no within-class association parameters). This slashes
    the number of parameters and hence the \emph{variance}; for large $p$ relative to $n$
    the variance saving usually outweighs the bias from ignoring dependence --- the
    bias--variance trade-off again.
  \end{solutionbox}
  \begin{solutionbox}[Solution 2(d) --- boundary shapes]
    \textbf{LDA:} linear boundaries (quadratic terms cancel due to the shared $\Sigma$).
    \textbf{QDA:} quadratic boundaries. \textbf{Naive Bayes:} generally non-linear but
    \emph{additive} boundaries --- log-odds of the form $\sum_j g_j(x_j)$ without
    interactions; only in special cases (e.g.\ Gaussian marginals with class-shared
    variances) does it become linear.
  \end{solutionbox}
\end{frame}

% ============================================================
% Problem 3
% ============================================================
\begin{frame}{Problem 3 --- restatement}
  \begin{exercise}[Problem 3 --- Confusion matrix and ROC curve (15 P)]
    Logistic-regression default classifier on $n=1{,}000$ customers with threshold
    $\hat{p}>0.5$:
    \begin{center}
    \begin{tabular}{llcc}
       & & \multicolumn{2}{c}{\textbf{True status}}\\
       & & Default & No default\\
      \midrule
      \textbf{Predicted} & Default & 60 & 45\\
       & No default & 40 & 855\\
      \midrule
       & Total & 100 & 900\\
    \end{tabular}
    \end{center}
    \begin{itemize}
      \item[(a)] TP/FP/FN/TN; accuracy and error rate; sensitivity; specificity;
        precision; FPR; interpret sensitivity and precision. \hfill (7 P)
      \item[(b)] Effect of lowering the threshold from 0.5 to 0.2 on sensitivity and
        specificity; business reason for accepting it. \hfill (4 P)
      \item[(c)] What does the ROC curve plot? Meaning of AUC $=0.5$ and AUC $=1$. \hfill (4 P)
    \end{itemize}
  \end{exercise}
\end{frame}

\begin{frame}{Problem 3 --- solution (a)}
  \begin{solutionbox}[Solution 3(a) --- the five metrics]
    Positive class ``default'': $\text{TP}=60$, $\text{FP}=45$, $\text{FN}=40$, $\text{TN}=855$.
    \begin{center}
    \begin{tabular}{llll}
      \toprule
      accuracy & $(60+855)/1000$ & $=0.915$ & error rate $=0.085$\\
      sensitivity (recall/TPR) & $60/(60+40)=60/100$ & $=0.600$ & \\
      specificity & $855/(855+45)=855/900$ & $=0.950$ & \\
      precision & $60/(60+45)=60/105$ & $=0.571$ & \\
      false positive rate & $45/(45+855)=45/900$ & $=0.050$ & $=1-\text{specificity}$\\
      \bottomrule
    \end{tabular}
    \end{center}
    \emph{Sensitivity:} $60\,\%$ of the true defaulters are detected --- $40\,\%$ are
    missed. \emph{Precision:} of all customers flagged as defaulters only $57.1\,\%$
    actually default; the remaining flags are false alarms.
  \end{solutionbox}
\end{frame}

\begin{frame}{Problem 3 --- solution (b) and (c)}
  \begin{solutionbox}[Solution 3(b) --- lowering the threshold to 0.2]
    More customers are flagged: true defaulters with $0.2<\hat{p}\le 0.5$ are now caught
    $\Rightarrow$ \textbf{sensitivity increases} (FN $\downarrow$); more non-defaulters
    are flagged too $\Rightarrow$ \textbf{specificity decreases} (FPR $\uparrow$).
    \emph{Business reason:} asymmetric costs --- a missed default (credit loss) is far
    more expensive than a false alarm (e.g.\ an unnecessary review); the expected cost can
    fall even if the raw error rate rises.
  \end{solutionbox}
  \begin{solutionbox}[Solution 3(c) --- ROC and AUC]
    The ROC curve plots the \emph{true positive rate} (sensitivity) against the
    \emph{false positive rate} ($1-$specificity) as the classification \emph{threshold}
    varies from 1 down to 0. $\text{AUC}=0.5$: diagonal --- no better than random
    guessing. $\text{AUC}=1$: perfect classifier --- some threshold attains sensitivity
    $=1$ and specificity $=1$ simultaneously. (AUC $=$ probability that a random defaulter
    is scored above a random non-defaulter.)
  \end{solutionbox}
\end{frame}

% ============================================================
% Problem 4
% ============================================================
\begin{frame}{Problem 4 --- restatement}
  \begin{exercise}[Problem 4 --- Resampling methods (18 P)]
    \begin{itemize}
      \item[(a)] $n=200$: number of model \emph{fits} for the validation-set approach
        (50/50 split) $\cdot$ 10-fold CV $\cdot$ LOOCV --- with justification. \hfill (4 P)
      \item[(b)] LOOCV vs.\ 10-fold CV: bias and variance; which is preferred in
        practice? \hfill (4 P)
      \item[(c)] Show $\Pr(i\in\text{bootstrap sample})=1-(1-1/n)^n$; evaluate at $n=5$;
        limit as $n\to\infty$. \hfill (5 P)
      \item[(d)] $B=3$ bootstrap estimates $\hat{\alpha}^{*}=0.4$ / $0.5$ / $0.6$
        (toy example; in practice $B\approx 1{,}000$): compute
        $\widehat{\operatorname{SE}}_B(\hat{\alpha})$; what does it estimate; why is the
        bootstrap useful? \hfill (5 P)
    \end{itemize}
  \end{exercise}
\end{frame}

\begin{frame}{Problem 4 --- solution (a) and (b)}
  \begin{solutionbox}[Solution 4(a) --- number of fits at n = 200]
    \begin{itemize}
      \item \textbf{Validation set: 1 fit} --- train once on the 100 training
        observations, evaluate on the 100 held-out ones.
      \item \textbf{10-fold CV: 10 fits} --- each fold of 20 is held out once; refit on
        the remaining 180 each time.
      \item \textbf{LOOCV: 200 fits} --- each observation is held out once; refit on the
        other 199 ($=n$-fold CV; for least squares a shortcut formula exists).
    \end{itemize}
  \end{solutionbox}
  \begin{solutionbox}[Solution 4(b) --- bias and variance of CV]
    \emph{Bias:} LOOCV trains on $n-1$ observations $\Rightarrow$ nearly unbiased;
    10-fold trains on $0.9\,n$ $\Rightarrow$ slightly overestimates the test error.
    \emph{Variance:} LOOCV is \emph{higher}: its $n$ fits use almost identical training
    sets, so the fold errors are strongly positively correlated --- averaging highly
    correlated quantities reduces variance little; the 10 folds overlap less.
    \emph{Practice:} $k=5$ or $10$ --- good bias--variance compromise and far cheaper.
  \end{solutionbox}
\end{frame}

\begin{frame}{Problem 4 --- solution (c) and (d)}
  \begin{solutionbox}[Solution 4(c) --- bootstrap inclusion probability]
    Each of the $n$ draws (with replacement) misses observation $i$ with probability
    $1-1/n$; draws are independent:
    $\Pr(i\notin\text{sample})=(1-1/n)^n \Rightarrow \Pr(i\in\text{sample})=1-(1-1/n)^n$.
    \[
    n=5:\ 1-(4/5)^5=1-0.32768=\mathbf{0.672};
    \qquad
    n\to\infty:\ 1-e^{-1}\approx 1-0.368=\mathbf{0.632}.
    \]
    A bootstrap sample contains roughly $2/3$ of the distinct original observations.
  \end{solutionbox}
  \begin{solutionbox}[Solution 4(d) --- bootstrap standard error]
    $\bar{\alpha}^{*}=(0.4+0.5+0.6)/3=0.5$;
    \[
    \widehat{\operatorname{SE}}_B(\hat{\alpha})
    =\sqrt{\tfrac{1}{3-1}\bigl[(-0.1)^2+0^2+(0.1)^2\bigr]}
    =\sqrt{0.02/2}=\sqrt{0.01}=\mathbf{0.100}.
    \]
    It estimates how much $\hat{\alpha}$ would vary over repeated samples from the
    population. The bootstrap needs no analytic formula and no new data --- it resamples
    the observed data and works for almost any statistic.
  \end{solutionbox}
\end{frame}

% ============================================================
% Problem 5
% ============================================================
\begin{frame}{Problem 5 --- restatement}
  \begin{exercise}[Problem 5 --- Model selection with Cp and BIC (12 P)]
    Best subset selection with $n=100$ and $\hat{\sigma}^2=4$ (from the full model):
    \begin{center}
    \begin{tabular}{lcccc}
      \toprule
      $d$ & 1 & 2 & 3 & 4\\
      \midrule
      $\operatorname{RSS}_d$ & 500 & 444 & 428 & 424\\
      \bottomrule
    \end{tabular}
    \end{center}
    \[
    C_p=\tfrac{1}{n}\bigl(\operatorname{RSS}+2d\hat{\sigma}^2\bigr),\qquad
    \operatorname{BIC}=\tfrac{1}{n}\bigl(\operatorname{RSS}+\log(n)\,d\,\hat{\sigma}^2\bigr),
    \qquad \log(100)\approx 4.605 .
    \]
    \begin{itemize}
      \item[(a)] $C_p$ for all four models. \hfill (4 P)
      \item[(b)] BIC for all four models. \hfill (4 P)
      \item[(c)] Which $d$ does each criterion select? Explain the disagreement. \hfill (2 P)
      \item[(d)] Why can training RSS alone never select $d$? \hfill (2 P)
    \end{itemize}
  \end{exercise}
\end{frame}

\begin{frame}{Problem 5 --- solution (a) and (b)}
  \begin{solutionbox}[Solution 5(a) --- Cp (penalty 8 per predictor)]
    $C_p=(\operatorname{RSS}_d+2d\hat{\sigma}^2)/100=(\operatorname{RSS}_d+8d)/100$:
    \begin{center}
    \begin{tabular}{lcccc}
      \toprule
      $d$ & 1 & 2 & 3 & 4\\
      \midrule
      $\operatorname{RSS}_d+8d$ & 508 & 460 & 452 & 456\\
      $C_p$ & 5.080 & 4.600 & \textbf{4.520} & 4.560\\
      \bottomrule
    \end{tabular}
    \end{center}
  \end{solutionbox}
  \begin{solutionbox}[Solution 5(b) --- BIC (penalty 18.421 per predictor)]
    $\operatorname{BIC}=(\operatorname{RSS}_d+4.605\cdot 4\cdot d)/100
      =(\operatorname{RSS}_d+18.421\,d)/100$:
    \begin{center}
    \begin{tabular}{lcccc}
      \toprule
      $d$ & 1 & 2 & 3 & 4\\
      \midrule
      $\operatorname{RSS}_d+18.421\,d$ & 518.421 & 480.841 & 483.262 & 497.683\\
      BIC & 5.184 & \textbf{4.808} & 4.833 & 4.977\\
      \bottomrule
    \end{tabular}
    \end{center}
  \end{solutionbox}
\end{frame}

\begin{frame}{Problem 5 --- solution (c) and (d)}
  \begin{solutionbox}[Solution 5(c) --- selected models and why they differ]
    $C_p$ picks $\mathbf{d=3}$ (4.520); BIC picks $\mathbf{d=2}$ (4.808).
    Both are RSS $+$ size penalty, but BIC charges $\log(n)\hat{\sigma}^2=18.421$ per
    predictor vs.\ $2\hat{\sigma}^2=8$ for $C_p$; since $\log(100)=4.605>2$, BIC
    penalises extra predictors more heavily and selects smaller models. Here the RSS drop
    from $d=2$ to $d=3$ is $16$: larger than $8$ (so $C_p$ accepts the third predictor)
    but smaller than $18.421$ (so BIC rejects it).
  \end{solutionbox}
  \begin{solutionbox}[Solution 5(d) --- why RSS alone fails]
    Adding a predictor can never increase the training RSS: RSS decreases
    \emph{monotonically} in $d$ (training $R^2$ increases monotonically), so RSS-based
    selection would always choose the largest model $d=4$ --- training error is an ever
    more optimistic estimate of test error as flexibility grows. Hence: penalise model
    size ($C_p$ / AIC / BIC / adjusted $R^2$) or estimate test error directly
    (validation set / cross-validation).
  \end{solutionbox}
\end{frame}

% ============================================================
% Problem 6
% ============================================================
\begin{frame}{Problem 6 --- restatement}
  \begin{exercise}[Problem 6 --- Ridge regression and the lasso (15 P)]
    Linear model with standardised predictors $X_1,\dots,X_p$.
    \begin{itemize}
      \item[(a)] Objectives of ridge and lasso; which coefficient is not penalised? \hfill (4 P)
      \item[(b)] Behaviour as $\lambda\to 0$ and $\lambda\to\infty$; which Chapter-3 method
        is recovered as $\lambda\to 0$? \hfill (3 P)
      \item[(c)] Why can the lasso set coefficients exactly to zero (geometry for $p=2$)? \hfill (3 P)
      \item[(d)] Why standardise the predictors first? \hfill (2 P)
      \item[(e)] Ridge fits at $\lambda=0.1$ and $\lambda=100$ (standardised coefficients)
        and a 10-fold CV grid: which fit has lower bias / lower variance; which $\lambda$
        should be chosen? \hfill (3 P)
    \end{itemize}
    \vspace{-1mm}
    \begin{center}
    \scriptsize
    \begin{tabular}{lrr}
      \toprule
       & $\lambda=0.1$ & $\lambda=100$\\
      \midrule
      $\hat{\beta}_1$ & $0.92$ & $0.31$\\
      $\hat{\beta}_2$ & $0.85$ & $0.29$\\
      $\hat{\beta}_3$ & $-0.41$ & $-0.12$\\
      $\hat{\beta}_4$ & $0.20$ & $0.06$\\
      \bottomrule
    \end{tabular}
    \hspace{10mm}
    \begin{tabular}{lcccc}
      \toprule
      $\lambda$ & 0.1 & 1 & 10 & 100\\
      \midrule
      CV error & 12.8 & 11.9 & 10.6 & 13.4\\
      \bottomrule
    \end{tabular}
    \end{center}
  \end{exercise}
\end{frame}

\begin{frame}{Problem 6 --- solution (a) and (b)}
  \begin{solutionbox}[Solution 6(a) --- the two objectives]
    \[
    \text{Ridge: }\min_{\beta}\ \operatorname{RSS}+\lambda\sum_{j=1}^{p}\beta_j^{2},
    \qquad
    \text{Lasso: }\min_{\beta}\ \operatorname{RSS}+\lambda\sum_{j=1}^{p}|\beta_j|,
    \]
    with $\operatorname{RSS}=\sum_{i}\bigl(y_i-\beta_0-\sum_j\beta_jx_{ij}\bigr)^2$ and
    tuning parameter $\lambda\ge 0$. The intercept $\beta_0$ is \emph{not} penalised.
  \end{solutionbox}
  \begin{solutionbox}[Solution 6(b) --- limiting behaviour]
    $\lambda\to 0$: the penalty vanishes; both approach the \emph{ordinary least-squares}
    fit of Chapter~3. $\lambda\to\infty$: the penalty dominates; all penalised
    coefficients go to zero (null model with intercept only). Ridge shrinks
    \emph{towards} zero (exactly zero only in the limit); the lasso sets more and more
    coefficients \emph{exactly} to zero as $\lambda$ grows --- flexibility decreases
    continuously along the $\lambda$ path.
  \end{solutionbox}
\end{frame}

\begin{frame}{Problem 6 --- solution (c) and (d)}
  \begin{solutionbox}[Solution 6(c) --- why the lasso yields exact zeros]
    Equivalent constrained form: minimise RSS subject to $|\beta_1|+|\beta_2|\le s$
    (lasso: a \emph{diamond} with corners on the axes) or $\beta_1^2+\beta_2^2\le s$
    (ridge: a \emph{disk} with no corners). The solution lies where the elliptical RSS
    contours around the OLS estimate first touch the region. Ellipses typically hit the
    diamond at a \emph{corner} --- where one coefficient is exactly zero (sparsity /
    variable selection). The disk has no corners, so the contact point generically has all
    coordinates non-zero: ridge shrinks but keeps every predictor.
  \end{solutionbox}
  \begin{solutionbox}[Solution 6(d) --- why standardise]
    OLS is scale-equivariant --- rescaling $X_j$ just rescales $\hat{\beta}_j$. The
    penalties depend on the coefficient \emph{magnitudes}: a predictor in small units has
    a large coefficient and would be shrunk more strongly for purely cosmetic reasons.
    Standardising to standard deviation one makes the penalty act comparably on all
    predictors and the fit independent of the units.
  \end{solutionbox}
\end{frame}

\begin{frame}{Problem 6 --- solution (e)}
  \begin{solutionbox}[Solution 6(e) --- bias-variance and the choice of lambda]
    $\lambda=0.1$: almost no shrinkage --- close to the OLS fit --- more flexible
    $\Rightarrow$ \textbf{lower bias and higher variance}. $\lambda=100$: strong shrinkage
    (e.g.\ $0.92\to 0.31$; each coefficient about a third of its size) $\Rightarrow$
    \textbf{higher bias and lower variance}. This is the Chapter-2 bias--variance
    trade-off controlled by the single knob $\lambda$.
    From the CV table the estimated test error is smallest at
    $\lambda=\mathbf{10}$ ($10.6<11.9<12.8<13.4$) --- choose $\lambda=10$; neither of the
    two fits in the coefficient table is optimal.
  \end{solutionbox}
  \begin{takeaway}[Big picture for the exam]
    Regularisation trades a little bias for a large variance reduction; $\lambda$ is
    always chosen by cross-validation --- never by looking at training RSS.
  \end{takeaway}
\end{frame}

\end{document}
